\documentclass{sprz}
\geometry{margin=25mm}
\begin{document}
\chapter*{Raport z laboratorium CI}
Technologie DevOps, laboratorium 02. Student: s30576.

\section*{Zakres sprawdzenia}
Mission Control korzysta z .NET 10, Angulara 21 i MySQL.
Pełny run CI wykonuje 19 gotowych testów HTTP API oraz 11 testów interfejsu.
Bazę uruchamia Testcontainers, a frontend korzysta z symulowanych odpowiedzi HTTP.

\section*{Pomiary cache}
Porównano ręczne uruchomienia na tej samej gałęzi. Klucze zawierały system,
architekturę, rzeczywistą wersję runtime i hash lockfiles.

\begin{center}
\begin{tabular}{lrrl}
\hline
Próba & npm ci [s] & NuGet [s] & Exact hit \\
\hline
Cold & 9 & 6 & nie \\
Warm & 7 & 4 & tak \\
Nowe lockfiles & 9 & 6 & nie \\
Nowy namespace & 7 & 6 & nie \\
\hline
\end{tabular}
\end{center}
Przy trafieniu cache nadal wykonano instalację, build i testy.
Zmiana lockfiles oraz namespace spowodowała użycie nowych kluczy.

\section*{Gate i metadane}
Check \texttt{ci-required} wymaga sukcesu detekcji zmian. Dla flagi
\texttt{true} dopuszcza wyłącznie wynik \texttt{success}, a dla
\texttt{false} wynik \texttt{skipped}.
Regresja walidacji nazwy misji spowodowała dwa nieudane testy i zablokowała PR.
Błąd detektora, brak flagi i nieoczekiwane pominięcie komponentu również
zostały odrzucone.

Osobny konsument potwierdził SHA, ID i attempt udanego producenta z main.
PR oraz anulowany run zostały pominięte, a test walidatora odrzucił schemat
w wersji 2. Linki do uruchomień są zapisane w dzienniku repozytorium.
\end{document}
